\bookchapter{Kapitel 11: Fotografie und Bildverwaltung}

\chapterintro{Ganz ehrlich: Die meisten Menschen fotografieren mit dem Smartphone weit unter dessen Potenzial. Nicht, weil Ihre Technik schlecht ist, sondern weil viele kleine Einstellungen und Gewohnheiten übersehen werden. Dieses Kapitel behandelt nichts Theoretisches, sondern nur Dinge, die Sie sofort anwenden können.}

\booksection{Haben Sie auch diese Fehler gemacht?}

Zuerst die häufigsten Probleme – schauen Sie, wie viele auf Sie zutreffen:

\begin{itemize}
  \item \textbf{Das Objektiv ist schmutzig}: Das ist wirklich das häufigste Problem. Beim Herausnehmen des Smartphones landen Fingerabdrücke auf dem Objektiv, die Aufnahmen wirken verschwommen, und man denkt, das Smartphone sei kaputt. Gewöhnen Sie sich an: Vor dem Fotografieren das Objektiv mit dem Ärmel abwischen.
  \item \textbf{Belichtung nie anpassen}: Im Standardmodus ist ein Tippen auf den Bildschirm die Fokussierung. Wenn Sie danach nach oben oder unten wischen, können Sie die Helligkeit einstellen. Viele Fotos sind zu dunkel oder zu hell, weil dieser Trick nicht bekannt ist.
  \item \textbf{Blitzlicht immer an}: Das LED-Blitzlicht des Smartphones liefert wirklich hässliche Ergebnisse – hartes, kaltes Licht, das Menschen gespenstisch erscheinen lässt. Außer bei völliger Dunkelheit: ausschalten. Bei schwachem Licht ist der Nachtmodus dem Blitzlicht weit überlegen.
  \item \textbf{Zoom mit zwei Fingern}: Das Auseinanderziehen mit zwei Fingern auf dem Bildschirm ist ein digitaler Zoom – im Grunde ein Ausschnitt des Bildes, der vergrößert wird, mit deutlichem Qualitätsverlust. Für entfernte Motive die Zoom-Tasten im Sucher verwenden (0,5x/1x/2x/3x) – das ist optischer Zoom ohne Qualitätsverlust.
\end{itemize}

\booksection{Kleine Tricks, die die Trefferquote sofort erhöhen}

Diese erfordern keine Theoriekenntnisse – einmal angewendet, sehen Sie den Effekt:

\begin{itemize}
  \item \textbf{Drittel-Regel / Gitterlinien}: Aktiviere in den Kameraeinstellungen die „Gitterlinien“ (auch Referenzlinien genannt). Platziere beim Fotografieren das Hauptmotiv auf den Schnittpunkten der vier Linien – das Bild wirkt sofort deutlich „richtiger“. Das ist die einfachste und wirksamste Regel der Fotografie, keine Mystik – sie entspricht nur den Sehgewohnheiten des menschlichen Auges.
  \item \textbf{Führen Sie mit Linien}: Achten Sie bei Landschafts- oder Straßenaufnahmen auf natürliche Linien im Bild – Straßen, Geländer, Wände, Baumreihen –, die vom Bildrand zum Hauptmotiv führen. Mit diesem Trick wirken Ihre Fotos, als hätten Sie Fotografie gelernt.
  \item \textbf{Nicht immer stehend fotografieren}: Die meisten fotografieren immer aus Augen höhe, deshalb wirken die Perspektiven eintönig. Probieren Sie, aus niedriger Perspektive zu fotografieren oder eine Erhöhung zu suchen und nach unten zu fotografieren – dieselbe Stelle wirkt völlig anders. Bei Kindern und Tieren ist das Herabknien besonders wichtig, um auf Augenhöhe zu sein.
  \item \textbf{Licht ist das wichtigste Element}: Mittags bei praller Sonne entstehen harte, unschöne Schatten im Gesicht. Das beste Licht ist in der „goldenen Stunde" am Morgen und Abend – das Licht ist schräg, die Farbe warm, alles sieht gut aus. Bewölktes Licht ist zwar nicht spektakulär, aber gleichmäßig und weich – für Porträts oft besser als grelles Sonnenlicht.
  \item \textbf{Nachtmodus richtig verwenden}: Im Nachtmodus macht das Smartphone automatisch mehrere Aufnahmen und fügt sie zusammen. Jetzt ist es entscheidend: Das Smartphone darf nicht zittern. An eine Wand lehnen, auf einen Tisch legen oder den Atem anhalten. Der Nachtmodus erfordert 2–3 Sekunden Ruhe – zittern macht das Bild unscharf.
\end{itemize}

\booksection{Porträtmodus und Pro-Modus}

\begin{itemize}
  \item \textbf{Porträtmodus}: Das Zweifinger-Pinchen auf dem Bildschirm ist digitaler Zoom: Das Bild wird zugeschnitten und vergroessert, die Qualitaet nimmt deutlich ab. Fuer entfernte Motive die Zoom-Tasten im Sucher verwenden (0.5x/1x/2x/3x). Dabei sind 0.5x (Ultraweitwinkel) und 1x (Hauptkamera) originale optische Linsen ohne Qualitaetsverlust; 2x/3x sind meist „optikqualitaets Zoom" – durch Cropping des hochaufloesenden Sensors realisiert, nahe an optischer Qualitaet, aber ohne eigene Linse. Ueber 3x ist es in der Regel Hybrid- oder Digitalzoom mit sichtbarem Qualitaetsverlust.
  \item \textbf{Pro-Modus}: Wenn Sie die volle Kontrolle über das Bild möchten, probieren Sie den Pro-Modus. Einige wichtige Parameter: Je niedriger die ISO, desto sauberer das Bild, aber es wird dunkler; je langsamer der Verschluss, desto mehr Licht, aber Verwackeln droht; der Weißabgleich steuert die Farbtemperatur – bei falscher Einstellung wird das Bild zu blau oder zu gelb. Die meisten brauchen den Pro-Modus nicht dauerhaft, aber die Kenntnis dieser Parameter hilft zu verstehen, was die Kamera tut.
\end{itemize}

\booksection{Nachbearbeitung: Nicht übertreiben, das rechte Maß ist richtig}

Viele haben ein Missverständnis gegenüber der Bildbearbeitung und finden, bearbeitet sei „nicht mehr echt“. Tatsächlich sieht ein Kamerabild ohnehin anders aus als das, was das Auge sieht – eine angemessene Anpassung dient dazu, das wiederherzustellen, was Sie damals gesehen haben, nicht etwas Künstliches zu erfinden.

\begin{itemize}
  \item \textbf{Snapseed}: Googles kostenloses Bildbearbeitungs-Tool, professionell, aber ohne Werbung. Mein täglicher Workflow: Erst mit „Bild anpassen" Helligkeit und Kontrast ziehen mit „Details" schärfen mit „Weißabgleich" Farbstiche korrigieren zum Schluss „Zuschneiden" auf die gewünschte Komposition. Die „Lokal"-Funktion von Snapseed ist sehr nützlich – Sie können nur einen Bereich des Bildes anpassen, ohne das gesamte Bild zu verändern.
  \item \textbf{Xingtu / Meitu XiuXiu}: Diese beiden eignen sich für Porträts. Aber ein Grundsatz: Schönheitsfilter nicht über 30 \% ziehen – darüber wird es künstlich, die Haut wirkt wie Plastik. Natürlich sieht besser aus als perfekt.
  \item \textbf{Lightroom Mobile}: Wenn Sie an Farben hohe Ansprüche stellen, ist das die professionellste Wahl. Kurvenkorrektur, HSL-Farbkanalanpassung und weitere Profifunktionen sind vorhanden. Für die meisten ist es aber etwas überdimensioniert.
  \item \textbf{Workflow-Empfehlung}: Meine Gewohnheit ist, erst gute Originalaufnahmen zu machen und zu Hause am Computer- oder iPad-Bildschirm in Ruhe zu bearbeiten – der Smartphone-Bildschirm ist zu klein und die Farben nicht genau genug. Wenn Sie unbedingt auf dem Smartphone bearbeiten müssen, stellen Sie die Bildschirmhelligkeit auf mittel bis hoch und bearbeiten Sie nicht im Energiesparmodus.
\end{itemize}

\booksection{Zu viele Fotos – was tun? Ordnungsstrategien}

Tausende Fotos in der Smartphone-Galerie, und man sucht ein bestimmtes ewig – dieses Erlebnis kennt jeder. Einige praktische Ordnungsmethoden:

\begin{itemize}
  \item \textbf{Suche nutzen}: iPhone und die meisten Android-Smartphones unterstützen die intelligente Suche in der Galerie. Suchen Sie mit Begriffen wie „Katze", „Strand", „Ausweis" – das System erkennt mit AI den Bildinhalt und findet die passenden Fotos. Auch nach Datum und Ort lässt sich suchen, z. B. „letzter Mai", „Peking". Diese Funktion kennen wenige, aber sie ist wirklich nützlich.
  \item \textbf{Regelmäßig Müll entfernen}: Screenshots, doppelte Fotos, unscharfe Ausschuss – das sind die wahren Platzfresser. Die iOS-Galerie hat eine Funktion „Duplikate automatisch zusammenführen"; bei Android suchen Sie in der Dateiverwaltung nach dem Ordner „screenshots" und löschen nicht benötigte Screenshots. Fünf Minuten pro Monat reichen, um mehrere GB freizugeben.
  \item \textbf{Alben erstellen, nicht nur auf Standardsortierung verlassen}: Die Standardsortierung nach Zeit ist nicht falsch, aber für zeitübergreifende Inhalte wie „Wachstum des Kindes", „Reisefotos" oder „wichtige Dokumente" sollten Sie manuell Alben anlegen. Apples „Gemeinsames Album" ermöglicht es der Familie, gemeinsam Fotos hinzuzufügen – eine Person erstellt es, alle Familienmitglieder können es sehen.
  \item \textbf{Backup ist die Grundlinie}: Fotos sind die unersetzlichsten Daten – Dateien können neu erstellt werden, Chatverläufe sind weg, wenn sie weg sind, aber Fotos der Kinder oder Videos der Älteren sind für immer verloren. Cloud-Backup ist Pflicht. iCloud-Fotos ist am sorgenfreiesten; bei Android hat jede Marke ihren eigenen Cloud-Dienst, oder Sie nutzen Baidu Netdisk / Alibaba Cloud Drive für automatische Backups. Der Grundsatz: Mindestens ein Cloud-Backup – nicht alle Eier in einen Korb legen.
\end{itemize}

\booksection{Screenshot und Bildschirmaufnahme}

\begin{itemize}
  \item \textbf{Screenshot}: Gleichzeitiges Drücken von Ein-/Aus-Taste und Lauter-Taste – universelle Methode für alle Smartphones. Viele Android-Smartphones unterstützen auch das Wischen mit drei Fingern nach unten oder das Doppeltippen mit dem Fingergelenk. Lange Screenshots (Scroll-Screenshots) sind auf Android ausgereift – nach dem ersten Screenshot auf die Schaltfläche „Scrollen" tippen, das Smartphone scrollt automatisch und fügt zu einem langen Bild zusammen, ideal für Chatverläufe und Webseiten.
  \item \textbf{Bildschirmaufnahme}: Fügen Sie im Kontrollzentrum „Bildschirmaufnahme" hinzu, tippen Sie zum Starten und zum Stoppen. Vor der Aufnahme langes Drücken der Aufnahmetaste, um auszuwählen, ob das Mikrofon Ton aufzeichnet – für ein Tutorial Mikrofon einschalten, für reine Bedienungsaufnahmen ausschalten. Das fertige Video wird automatisch in der Galerie gespeichert.
\end{itemize}

\begin{tcbpractice}{Üben Sie: Ein vollständiger Durchlauf aus Fotografieren + Bearbeiten + Ordnen}
\begin{enumerate}
  \item Aktivieren Sie die „Gitterlinien" in den Kameraeinstellungen und machen Sie an einem Ort mit natürlichem Licht drei Fotos (Objektiv wischen, Belichtung anpassen).
  \item Wählen Sie das beste Foto und nehmen Sie mit Snapseed Grundeinstellungen vor (Helligkeit, Kontrast, Zuschneiden).
  \item Erstellen Sie für dieses Foto ein eigenes Album mit dem Namen „Meine Werke".
  \item Prüfen Sie, ob das Cloud-Backup aktiviert ist, und stellen Sie sicher, dass das Foto in der Cloud gespeichert ist.
\end{enumerate}

\end{tcbpractice}
